\documentclass[profile=editorial]{atelier}
\AtelierUseACGNBibliography
% 防退化说明：仓库编译测试与公开命令都从根目录执行；书目文件使用根目录
% 相对路径，确保示例在脚本和手工命令中读取同一份数据。
\addbibresource{examples/gbt7714-acgn-profile/references.bib}

\title{ATX-ACGN-REF 0.4}
\subtitle{GB/T 7714-2025 + ACGN Media Index · Real-world Citation Showcase}
\author{AtelierTeX Project}
\publicationdate{2026-08-18}
\versionnote{Public real-world fixture · first-party and authoritative sources preferred}

\begin{document}
\maketitle

\section{正文引用方式}

普通作品、角色档案和剧情材料可以直接使用数字顺序制引用。游戏本体使用
\verb|\cite{gakumas_game}|，例如《学園アイドルマスター》\cite{gakumas_game}；
角色档案使用 \verb|\cite{hiro_character_file}|，例如篠澤広官方档案\cite{hiro_character_file}。

作品内部位置使用 BibLaTeX 的 postnote 给出 Locator。例如：

\begin{verbatim}
\cite[STEP1 / Episode 8]{hiro_commu_step1_08}
\cite[Episode 1]{umamusume_anime_ep01}
\cite[Vol. 1]{haruhi_manga_vol1}
\end{verbatim}

实际渲染示例：剧情节点\cite[STEP1 / Episode 8]{hiro_commu_step1_08}，
动画单集\cite[Episode 1]{umamusume_anime_ep01}，漫画单行本\cite[Vol. 1]{haruhi_manga_vol1}。
同一论断引用多项材料时可使用 \verb|\cite{key1,key2}|，例如官方曲目页面与实体单曲页面\cite{hiro_song_koukei,hiro_1st_single}。

\section{ACGN 与跨媒介来源}

本样例使用真实官方页面演示游戏、角色档案、游戏内剧情、PV、设定资料、虚构校园地图、互动活动、歌曲、单曲、MV、新闻、Live、音频文件、Web 漫画索引、动画单集、漫画单行本、轻小说和开发者访谈。对应来源包括：
学园偶像大师官方站\cite{hiro_character_pv,hatsuboshi_schoolguide,hatsuboshi_campus_map,hatsuboshi_entrance_exam,gakumas_comics_index}、
学园偶像大师 Label 与 Live 官方页\cite{gakumas_mv_news,gakumas_first_single_news,gakumas_debut_live_hatsukoi,gakumas_audio_instrumental}、
赛马娘动画官方 Story\cite{umamusume_anime_ep01}、KADOKAWA 书目页\cite{haruhi_manga_vol1,haruhi_novel_vol1}和 Nintendo 开发者访谈\cite{mario_wonder_developer_interview}。

\section{普通学术文献与 GB/T 类型}

普通期刊、图书、会议论文、学位论文和报告与 ACGN 来源共用同一参考文献表。这里分别演示期刊论文\cite{mnih2015dqn}、图书\cite{sutton_barto_2018}、会议论文\cite{vaswani2017attention}、学位论文\cite{vanhorn2019visipedia}和技术报告\cite{nist_ai_rmf_2023}。

标准、数据集、专利、报纸和数据库使用各自的 BibLaTeX entry type：GB/T 7714-2025\cite{gbt7714_2025}、UCI Iris 数据集\cite{uci_iris}、PageRank 专利\cite{page2001pagerank}、中国科学报数学书评\cite{qiyun2024math}以及 dblp 数据库\cite{dblp_database}。

\section{数字文件与软件工程对象}

JSON、CSV、YAML、Office Open XML、PNG、SVG、WebVTT、ZIP 和 Jupyter Notebook 的公开规范或文档用于演示文件载体引用\cite{rfc8259_json,rfc4180_csv,yaml_122,ecma376_ooxml,w3c_png3,w3c_svg2,w3c_webvtt,pkware_zip_appnote,jupyter_nbformat}。
软件发行与代码仓库使用 `RELEASE / REPOSITORY` Media Tag，例如 Python 3.14.0\cite{python_3140}、AtelierTeX\cite{ateliertex_repo} 与 biblatex-gb7714-2025\cite{biblatex_gb7714_repo}。

\section{预印本与公共样例边界}

\ifstrequal{\AtelierActualGBStyle}{gb7714-2025}{%
GB/T 7714-2025 环境原生演示 \verb|@preprint|：Attention Is All You Need 的 arXiv 版本\cite{vaswani2017attention_preprint}。%
}{%
当前 TeX 环境使用 \texttt{\AtelierActualGBStyle} 兼容渲染；\texttt{@preprint} 记录保留在 \texttt{references.bib} 中，正式 2025 样式下进入参考文献表。%
}

`FAN LOCATOR` 保留一个格式 fixture，用于展示字段结构；公共样例不把未经项目人工复核的 fan wiki / fan database 当作证据来源。

\nocite{fan_locator_fixture}
\printbibliography

\end{document}
